\section{Introduction}
\subsection{Contexte et thématique}
Le présent mini-projet de séries temporelles porte sur l'analyse et la prévision
de trois séries horaires du système électrique allemand, mises à disposition
par la plateforme de transparence \textbf{ENTSO-E} :
\begin{itemize}
  \item \texttt{Load} : la consommation électrique nationale (charge nette) ;
  \item \texttt{Wind} : la production éolienne terrestre et en mer ;
  \item \texttt{Solar} : la production photovoltaïque.
\end{itemize}
Toutes les séries sont exprimées en \textbf{mégawatts (MW)} et couvrent la
période du \textbf{1\ier{} janvier 2015} au \textbf{30 septembre 2020}, soit
\(50\,400\) observations à pas horaire, horodatées en \textbf{UTC} (ce qui
élimine toute ambiguïté liée aux changements d'heure été/hiver).

Le choix de l'énergie se justifie par la richesse structurelle des données :
forte saisonnalité journalière et hebdomadaire (cycles d'activité humaine pour
la charge), saisonnalité annuelle et météorologique (éolien et solaire),
volatilité par grappes (épisodes venteux, journées ensoleillées) et
interdépendances économiques entre consommation et production renouvelable.
Ce contexte autorise un spectre complet de méthodes : modèles univariés
(ARIMA/SARIMA), modèles multivariés dépendant du temps (VAR/VECM),
modélisation de la volatilité (GARCH) et apprentissage profond (LSTM/GRU).

\subsection{Problématique}
La transition énergétique allemande (\emph{Energiewende}) fait de l'équilibre
offre/demande un enjeu critique. Trois questions guident le travail :
\begin{enumerate}
  \item Peut-on \textbf{prévoir} ces séries à court (\(h=10\) et \(h=24\)
        heures), moyen et très long terme (\(h=5\,040\) heures, soit environ
        \(7\) mois) ?
  \item Existe-t-il des \textbf{relations dynamiques} (causalité, cointégration)
        entre consommation et productions renouvelables ?
  \item Les \textbf{modèles profonds} (LSTM/GRU) surpassent-ils les modèles
        classiques, et de façon statistiquement significative ?
\end{enumerate}

\subsection{Objectifs}
Le travail poursuit quatre objectifs, alignés sur le cahier des charges :
\begin{enumerate}
  \item Réaliser le \textbf{socle classique en R} (sections 3.1 ARIMA/SARIMA,
        3.2 VAR / causalité de Granger / VECM, 3.3 ARCH/GARCH) ;
  \item \textbf{Transposer} l'intégralité de ce socle en \textbf{Python}, afin
        de vérifier la reproductibilité des résultats (\emph{double
        implémentation}) ;
  \item Étendre l'analyse aux \textbf{modèles profonds LSTM/GRU} en Python ;
  \item Conduire une \textbf{évaluation comparative} rigoureuse (métriques
        MSE/RMSE/MAE/MAPE, walk-forward à un pas, test de
        \textbf{Diebold--Mariano}) et restituer l'ensemble dans un rapport
        scientifique avec tableaux et figures issus des scripts.
\end{enumerate}

\subsection{Organisation du rapport}
La suite est organisée comme suit. La section~\ref{sec:methodo} détaille les
données, le prétraitement et l'ensemble des méthodologies. La
section~\ref{sec:resultats} présente les résultats empiriques (analyse
exploratoire, ARIMA/SARIMA, VAR/VECM, GARCH, LSTM/GRU). La
section~\ref{sec:evaluation} compare les modèles. La
section~\ref{sec:conclusion} conclut et discute les limites. Les annexes
recensent l'ensemble des figures et tableaux produits.
